\documentclass[../readings.tex]{subfiles}

\begin{document}

\subsection{Iterative Methods for Linear Systems}
\label{sub:iterative-basics}

Alternative to direct solvers for large, sparse systems where matrix factorization is impractical due to memory or compute limits.

\textBox{%
Direct methods factor the matrix. Iterative methods repeatedly improve an approximate solution using cheap operations, usually matrix-vector products. They are preferred when $\bA$ is large, sparse, or available only as a function that computes $\bA\bv$.
}

\subsubsection{Direct vs.\ Iterative Solvers}

\addRemark{%
\textbf{(Sparse fill-in)}\enspace
Even if $\bA$ is sparse, its LU or Cholesky factors are typically much denser (\textbf{fill-in}). For a 2D Poisson problem on an $n \times n$ grid:
\begin{itemize}
    \item $\bA$ has $O(n^2)$ nonzeros.
    \item LU factors have $O(n^3)$ nonzeros.
\end{itemize}
Iterative methods exploit sparsity by using only matrix-vector products ($O(\text{nnz})$).
}

\addRemark{%
\textbf{(Matrix-free viewpoint)}\enspace Many scientific codes never assemble $\bA$ explicitly. They only provide a routine that maps $\bv\mapsto\bA\bv$. Krylov methods such as CG and GMRES are designed for this setting.
}

\addDef{Krylov Subspace}{%
Given a matrix $\bA$ and a vector $\bv$, the $k$-th Krylov subspace is
\begin{align*}
    \mathcal{K}_k(\bA,\bv)
    =
    \operatorname{span}\{\bv,\bA\bv,\bA^2\bv,\dots,\bA^{k-1}\bv\}.
\end{align*}
It is the smallest subspace that contains $\bv$ and is built by repeated multiplication by $\bA$.
}
\label{def:krylov-subspace-general}

\addThm{Krylov Approximation Principle}{%
Let $\br^{(0)}=\bb-\bA\bx^{(0)}$. Any iterate of the form
\begin{align*}
    \bx^{(k)}=\bx^{(0)}+p_{k-1}(\bA)\br^{(0)},
\end{align*}
where $p_{k-1}$ is a polynomial of degree at most $k-1$, lies in the affine space
\begin{align*}
    \bx^{(0)}+\mathcal{K}_k(\bA,\br^{(0)}).
\end{align*}
Thus a matrix-free method can improve $\bx^{(0)}$ by choosing increasingly good corrections from Krylov subspaces using only matrix-vector products.
}
\label{thm:krylov-approximation-principle}

\addExcr{%
\begin{enumerate}
    \item Write $p_{k-1}(t)=c_0+c_1t+\cdots+c_{k-1}t^{k-1}$ and expand $p_{k-1}(\bA)\br^{(0)}$.
    \item Use Def~\ref{def:krylov-subspace-general} to show that this correction belongs to $\mathcal{K}_k(\bA,\br^{(0)})$.
    \item List $\mathcal{K}_1$, $\mathcal{K}_2$, and $\mathcal{K}_3$ explicitly. What new vector must be computed to move from $\mathcal{K}_k$ to $\mathcal{K}_{k+1}$?
    \item Explain why Thm~\ref{thm:krylov-approximation-principle} makes Krylov methods compatible with the matrix-free viewpoint above.
    \item Preview the later chapters: CG chooses the best correction in an energy norm for SPD systems, while GMRES chooses the correction with smallest residual norm.
\end{enumerate}
}
\label{ex:krylov-approximation-principle}

\addDef{Stationary Iteration}{%
An iteration of the form \textbf{$\bx^{(k+1)} = \bT\bx^{(k)} + \bc$}, where:
\begin{itemize}
    \item $\bT = \bI - \bM^{-1}\bA$ is the iteration matrix.
    \item $\bc = \bM^{-1}\bb$ is the updated right-hand side.
    \item $\bM$ is the \textbf{splitting matrix} (ideally $\bM \approx \bA$ and $\bM^{-1}$ is cheap).
\end{itemize}
}
\label{def:stationary-iteration}

\textBox{%
Stationary iterations are fixed-point iterations. Instead of solving $\bA\bx=\bb$ directly, we rewrite the problem as $\bx=\bT\bx+\bc$ and repeatedly apply the map. The exact solution is a fixed point of this map.
}

\addRemark{%
\textbf{(Splitting form)}\enspace If $\bA=\bM-\mathbf{N}$, then
\begin{align*}
    \bA\bx=\bb
    \quad \Longleftrightarrow \quad
    \bM\bx=\mathbf{N}\bx+\bb.
\end{align*}
The iteration is
\begin{align*}
    \bx^{(k+1)}=\bM^{-1}\mathbf{N}\bx^{(k)}+\bM^{-1}\bb.
\end{align*}
The art is choosing $\bM$ easy to invert while making the iteration converge quickly.
}

\subsubsection{Convergence and Spectral Radius}

\addDef{Spectral Radius}{%
$\rho(\bB) = \max \{|\lambda| : \lambda \in \sigma(\bB)\}$.
}
\label{def:spectral-radius}

\addRemark{%
\textbf{(Spectral radius intuition)}\enspace The eigenvalue of largest magnitude controls the long-run behavior of repeated multiplication. If every eigenvalue of $\bT$ lies inside the unit disk, powers of $\bT$ decay and the iteration converges.
}

\addThm{Gelfand's Formula}{%
For any submultiplicative norm, $\rho(\bB) = \lim_{k\to\infty} \|\bB^k\|^{1/k}$.
For large $k$, $\|\bB^k\| \approx \rho(\bB)^k$.
}
\label{thm:gelfand}

\addThm{Fundamental Convergence Condition}{%
The iteration $\bx^{(k+1)} = \bT\bx^{(k)} + \bc$ converges for any $\bx^{(0)}$ iff \textbf{$\rho(\bT) < 1$}.
}
\label{thm:iteration-convergence}

\addPrf{%
Let $\bx^*$ be the exact solution satisfying $\bx^* = \bT\bx^* + \bc$. Subtracting this from the iteration formula gives the error evolution:
\begin{align*}
    \bx^{(k+1)} - \bx^* &= (\bT\bx^{(k)} + \bc) - (\bT\bx^* + \bc) \\
    \be^{(k+1)} &= \bT \be^{(k)}.
\end{align*}
Applying this relation recursively from $k=0$:
\begin{align*}
    \be^{(k)} = \bT^k \be^{(0)}.
\end{align*}
The error $\be^{(k)} \to \bzero$ for any $\be^{(0)}$ iff $\bT^k \to \bzero$ as $k \to \infty$, which is guaranteed iff $\rho(\bT) < 1$ by Gelfand's formula.
}

\addRemark{%
\textbf{(Iteration count)}\enspace
To reduce error by $10^{-p}$, the required iterations are $k \approx \frac{p}{\log_{10}(1/\rho(\bT))}$. As $\rho(\bT) \to 1^-$, the cost explodes.
}

\subsubsection{Stopping Criteria}

\textBox{%
Iterative methods need a stopping rule. The true error $\|\bx^{(k)}-\bx^*\|$ is usually unknown, so algorithms monitor the residual $\br^{(k)}=\bb-\bA\bx^{(k)}$.
}

\addRemark{%
\textbf{(Common stopping test)}\enspace Stop when
\begin{align*}
    \frac{\|\br^{(k)}\|}{\|\bb\|} \leq \texttt{tol}.
\end{align*}
This is a relative residual test. It is cheap and scale-aware, but it is not the same as a relative error test unless the problem is well-conditioned.
}

\addRemark{%
\textbf{(Preconditioning preview)}\enspace Most useful iterative solvers are preconditioned. A preconditioner $\bM\approx\bA$ should be cheap to solve with and should make the transformed system easier for the iterative method. Good preconditioning changes the problem from "possible but slow" to practical.
}

\addExcr{%
\begin{enumerate}
    \item For $\bA = \begin{pmatrix} 3 & 1 \\ 1 & 3 \end{pmatrix}$, find $\rho(\bT)$ for the splitting $\bM = 3\bI$.
    \item If $\rho(\bT) = 0.99$, how many iterations are needed to gain 4 digits of precision?
    \item Start $\bx^{(0)}=\bzero$ and solve $\bA\bx = (4, 4)^T$ via 3 steps of this iteration.
\end{enumerate}
}

\subsubsection{Fill-in and Sparse Storage}

\addExcr{%
\begin{enumerate}
    \item \textbf{Fill-in Demo:}\enspace Generate a large sparse tridiagonal matrix. Compare nonzeros in $\bA$ vs.\ LU factors.
    \item \textbf{Timing:}\enspace Use \texttt{scipy.sparse.linalg.spsolve} vs.\ dense \texttt{np.linalg.solve} for $n=2000$. Record memory and time.
\end{enumerate}
}

\end{document}
